% 5-5 Lesson Quiz, short form. Items 1, 4 and 5 of the TE's "5-5 Lesson Quiz: Function Operations" (2024 TE page 268A),
% the items the kept Examples 5 and 6 teach. Items 2 and 3 (Examples 2 and 3, cut) are listed on the key page for a Form B.
\documentclass[11pt]{article}
\usepackage[margin=0.7in]{geometry}
\usepackage{amsmath,amssymb}
\usepackage{enumitem}
\usepackage{tikz}
\usepackage{pgfplots}
\pgfplotsset{compat=1.18}
\usepackage{fancyhdr}
\usepackage{parskip}
\setlist[enumerate]{leftmargin=*, itemsep=6pt}
\pagestyle{fancy}
\fancyhf{}
\fancyfoot[L]{\small\textit{en}Vision\textsuperscript{\textregistered} Algebra 2}
\fancyfoot[C]{\small 5-5 Lesson Quiz, short form}
\fancyfoot[R]{\small\thepage}
\renewcommand{\headrulewidth}{0pt}
\newcommand{\choice}[1]{\tikz[baseline=-0.6ex]\node[draw, circle, inner sep=0.6pt, minimum size=3.3mm, font=\scriptsize\bfseries] {#1};\;}
\newcommand{\checkbox}{$\square$\;}

\begin{document}

\noindent\textbf{Name}\; \rule{6cm}{0.4pt} \hfill Period \rule{1.2cm}{0.4pt} \quad Date \rule{2.2cm}{0.4pt}

\vspace{4pt}\noindent\rule{\textwidth}{0.6pt}

\vspace{2pt}\noindent{\LARGE\textbf{5-5 \; Lesson Quiz: Function Operations}} \hfill {\small short form}

\vspace{4pt}\noindent\rule{\textwidth}{0.6pt}

\vspace{6pt}
\textbf{1.} Let $f(x)=3x^2-2x+6$ and $g(x)=7x-4$. Identify the rule for $g(f(x))+g(x)$.

\vspace{4pt}
\begin{minipage}[t]{0.48\textwidth}
\begin{itemize}[leftmargin=1.2em, itemsep=6pt]
\item[] \choice{A} $21x^2-7x+34$
\item[] \choice{B} $21x^2-21x+50$
\end{itemize}
\end{minipage}\hfill
\begin{minipage}[t]{0.48\textwidth}
\begin{itemize}[leftmargin=1.2em, itemsep=6pt]
\item[] \choice{C} $147x^2-182x+62$
\item[] \choice{D} $21x^3-26x^2+57x-28$
\end{itemize}
\end{minipage}

\vspace{4pt}
Show your work:

\vspace{2.6cm}

\noindent\rule{\textwidth}{0.4pt}

\vspace{6pt}
\textbf{4.} Let $h(x)=\sqrt{x+3}$ and $k(x)=2x+7$. Find the value $\dfrac{h(k(3))-k(h(1))}{k(h(-2))-h(k(-3))}$.

\vspace{4pt}
\hspace*{1.2em}\choice{A} $\dfrac{2\sqrt6-2\sqrt3}{\sqrt6-7}$ \hfill \choice{B} $4$ \hfill \choice{C} $2\sqrt6-\sqrt{12}+7$ \hfill \choice{D} $-1$\hspace*{2em}

\vspace{4pt}
Show your work:

\vspace{2.2cm}

\noindent\rule{\textwidth}{0.4pt}

\vspace{6pt}
\textbf{5.} Let $f(x)=3x^2+2$ and $g(x)=\sqrt{x-4}$. What is the rule for the composition $f\circ g$?

\vspace{4pt}
\begin{itemize}[leftmargin=1.2em, itemsep=2pt]
\item[] \choice{A} $(f\circ g)(x)=\sqrt{3x^2-2}$, domain is all real numbers
\item[] \choice{B} $(f\circ g)(x)=\sqrt{3x^2-2}$, domain is $x\geq4$
\item[] \choice{C} $(f\circ g)(x)=3x-10$, domain is all real numbers
\item[] \choice{D} $(f\circ g)(x)=3x-10$, domain is $x\geq4$
\end{itemize}

\vspace{4pt}
Show your work, and explain in one sentence how you found the domain:

\vspace{2.2cm}

\vfill
\noindent\textbf{Scoring:} Question 1 and Question 5 are 10 points each. Question 4 is a bonus of up to 5 points. Total out of 20.

\newpage
\noindent{\LARGE\textbf{5-5 Lesson Quiz, short form: answer key}}

\vspace{4pt}\noindent\rule{\textwidth}{0.6pt}

\begin{enumerate}
\item[\textbf{1.}] \textbf{A, $21x^2-7x+34$.} $g(f(x))=7(3x^2-2x+6)-4=21x^2-14x+42-4=21x^2-14x+38$; adding $g(x)=7x-4$ gives $21x^2-7x+34$. Choice C is $f(g(x))=3(7x-4)^2-2(7x-4)+6=147x^2-182x+62$ (the wrong order, and $g(x)$ never added); D is $f(x)\cdot g(x)+g(x)$ (a product, not a composition). Savvas quiz item 1, DOK 3, skills F52, F26; taught by Example 5B (substitute the whole inside rule) and Example 6 ($f\circ g$ versus $g\circ f$), with QQD (starred) for the sum. \textbf{10 points:} 6 for A, 4 for the work ($7(3x^2-2x+6)-4$ written and the like terms combined).
\item[\textbf{4.}] \textbf{D, $-1$.} Inside values: $k(3)=13$, $h(1)=2$, $h(-2)=1$, $k(-3)=1$. Then $h(13)=\sqrt{16}=4$, $k(2)=11$, $k(1)=9$, $h(1)=\sqrt4=2$, so the value is $\dfrac{4-11}{9-2}=\dfrac{-7}{7}=-1$. Savvas quiz item 4, DOK 3, skills F53, F26; taught by working inside out (Example 5A's check line, Example 6's test price, the packet's Exit question). \textbf{Bonus 5:} 5 for D with the inside and outside values shown; 3 for D with some of them shown; 2 for D with no work.
\item[\textbf{5.}] \textbf{D, $3x-10$, domain $x\geq4$.} $f(g(x))=3\left(\sqrt{x-4}\right)^2+2=3(x-4)+2=3x-10$. The simplified rule is linear, but the inside function $g(x)=\sqrt{x-4}$ is defined only for $x\geq4$, so the domain of $f\circ g$ is $x\geq4$ (choice C is the domain trap). Choices A and B put $f$ inside $g$: $g(f(x))=\sqrt{3x^2+2-4}=\sqrt{3x^2-2}$. Savvas quiz item 5, DOK 2, skills F64, F26; taught by Example 5A and the packet's domain checks 1 and 2 (in Example 5A the outside square root sets the limit; here the inside function has the square root, so $x$ must be in the domain of $g$). \textbf{10 points:} 6 for D, 4 for the work (the substitution written and the domain taken from the inside function).
\end{enumerate}

\textbf{Items left off this form} (available for a retake or Form B): Savvas item 2, the revenue function for demand $d(x)=500-2x$ and price $p(x)=x+5$, $R(x)=2{,}500+490x-2x^2$ (answer 2,500; 490; 2; DOK 3, skills F53, F26; Example 2, cut in the fast packet); Savvas item 3, the rule and domain of $\frac fg$ for $f(x)=-3x-6$ and $g(x)=x^2-x-6$ (Part A answer D, $\frac{-3}{x-3}$, DOK 3; Part B answer B and E, $3$ and $-2$ are not in the domain, DOK 2; skills F53, F26; Example 3, cut in the fast packet). The TE quiz has no graphs, so nothing was redrawn.

\textbf{IXL replacement:} SmartScore 70 or above on EKJ, RSP, P9T and QQD replaces this quiz with 20/20.

\end{document}
